\begin{textsample}{POS Dim 4 – human – Score 6.00 – rj/rj\_ef\_linguagens\_lp\_1\_p4.txt}  \label{ex:f4_pos_004}
Nesse sentido, através de um trabalho interdisciplinar no qual o professor pode atuar como mediador, a produção de conhecimento pelos estudantes se dará através dos usos das variadas linguagens presentes nos diferentes tipos \textbf{textuais} utilizados em sala de aula, \textbf{associadas} às demandas de \textbf{leitura} e \textbf{escrita} advindas da prática social que permeiam o ambiente em que vivem, dentro e fora dela.
A escola \textbf{precisa} ser um local de interlocução, ou seja, \textbf{precisa} \textbf{abrir} espaços para o uso de diferentes linguagens, interações e discursos, ouvindo o que professores e estudantes têm a dizer a partir de suas experiências e vivências, já que ela é o lugar de excelência para o aprendizado.
Sendo a língua viva, entendida como algo inacabado, em constante transformação e que acontece no discurso, é de suma importância levar em consideração seus diferentes usos, em diferentes contextos e situações.

% matched lemmas: abrir, associar, escrita, leitura, precisar, textual
\end{textsample}
